\chapter{Ogólne twierdzenie Stokesa}
\label{ma2:ch:stokes}

W rozdziale \ref{ma2:ch:differential-forms}
zdefiniowaliśmy pochodną zewnętrzną oraz całkę
formy po zorientowanej podrozmaitości.
Teraz połączymy te dwie operacje.
Ogólne twierdzenie Stokesa mówi, że całka
z pochodnej zewnętrznej po podrozmaitości
jest równa całce z pierwotnej formy po jej brzegu.
Twierdzenia Newtona--Leibniza, Greena, Gaussa
o dywergencji i klasyczne twierdzenie Stokesa
dla pola wektorowego okażą się jego przypadkami.

\section{Intuicja: znoszenie się wkładów wewnętrznych}

Jeżeli podzielimy prostokąt na dwa prostokąty
i zorientujemy oba zgodnie z orientacją płaszczyzny,
to wspólny bok otrzymuje dwie przeciwne orientacje.
W sumie całek po brzegach obu prostokątów wkłady
wspólnego boku się znoszą. Pozostaje całka
po zewnętrznym brzegu całego prostokąta.
Zjawisko to ilustruje mechanizm twierdzenia Stokesa.
Nie jest jednak jego ogólnym dowodem: dla dowolnej
podrozmaitości użyjemy parametryzacji i gładkich
funkcji o małych nośnikach.

\begin{figure}[htbp]
\centering
\begin{tikzpicture}[>=Latex,scale=1.1]
\fill[blue!12] (0,0) rectangle (2.5,2);
\fill[orange!15] (2.5,0) rectangle (5,2);
\draw[thick] (0,0) rectangle (5,2);
\draw[gray,dashed] (2.5,0)--(2.5,2);
\foreach \x in {0,2.5}{
 \draw[->,very thick,blue!70!black] (\x+.35,0)--(\x+2.15,0);
 \draw[->,very thick,blue!70!black] (\x+2.15,2)--(\x+.35,2);
}
\draw[->,very thick,blue!70!black] (0,1.65)--(0,.35);
\draw[->,very thick,blue!70!black] (5,.35)--(5,1.65);
\draw[->,very thick,red!75!black] (2.38,.35)--(2.38,1.65);
\draw[->,very thick,red!75!black] (2.62,1.65)--(2.62,.35);
\node at (1.25,1) {$D_1$};
\node at (3.75,1) {$D_2$};
\node[red!70!black,below] at (2.5,-.35) {przeciwne wkłady};
\end{tikzpicture}
\caption{Orientacje wspólnej krawędzi są przeciwne.
W sumie całek brzegowych pozostaje zewnętrzny brzeg.}
\label{ma2:fig:stokes-cancellation}
\end{figure}

\section{Wzór lokalny w przestrzeni euklidesowej}

W całym rozdziale przyjmujemy standardową orientację
$\R^k$, a na półprzestrzeni
$H^k=\{u\in\R^k:u_k\ge0\}$ orientację indukowaną
przez $\R^k$. Na jej brzegu stosujemy regułę
z \ref{ma2:prop:boundary-orientation}.
Dla $k=1$ całka po brzegu półprostej oznacza
wartość funkcji w zerze ze znakiem $-1$.

\begin{lemat}[Twierdzenie Stokesa w modelach lokalnych]
\label{ma2:lem:stokes-local}
Niech $k\ge1$ i niech $X$ będzie $\R^k$ albo $H^k$.
Jeżeli $\eta$ jest formą stopnia $k-1$ klasy $C^1$
o zwartym nośniku w $X$, to
\[
 \int_X\mathrm d\eta=\int_{\partial X}j^*\eta,
\]
gdzie $j:\partial X\to X$ jest odwzorowaniem
włączenia, a całka po zbiorze pustym wynosi zero.
\end{lemat}
\begin{proof}
Piszemy jednoznacznie
\[
 \eta=\sum_{i=1}^k(-1)^{i-1}a_i(u)
 \,\mathrm du_1\wedge\cdots\wedge
 \widehat{\mathrm du_i}\wedge\cdots\wedge\mathrm du_k.
\]
Dla $k=1$ pusty iloczyn zewnętrzny oznacza $1$.
Z definicji pochodnej zewnętrznej
\ref{ma2:def:exterior-derivative}
wynika
\[
 \mathrm d\eta=\left(\sum_{i=1}^k\partial_i a_i\right)
 \mathrm du_1\wedge\cdots\wedge\mathrm du_k.
\]
Istotnie, pochodne względem $u_j$ dla $j\ne i$
dają powtórzony czynnik $\mathrm du_j$ i znikają.
Przeniesienie $\mathrm du_i$ na $i$-tą pozycję
daje znak $(-1)^{i-1}$, który znosi znak
w definicji składnika $\eta$.

Współczynniki $a_i$ i ich pierwsze pochodne
mają zwarty nośnik, ponieważ znikają lokalnie
poza nośnikiem $\eta$. Są więc całkowalne.
W standardowej orientacji całka formy
najwyższego stopnia jest całką z jej współczynnika
względem $\lambda_k$, na mocy
\ref{ma2:thm:integral-top-form-chart}.

Dla $X=\R^k$ stosujemy twierdzenie Fubiniego
i twierdzenie Newtona--Leibniza w każdej zmiennej
$u_i$. Przy ustalonych pozostałych współrzędnych
funkcja $a_i$ znika dla dostatecznie dużego
$|u_i|$, więc całka z $\partial_i a_i$
po całej prostej wynosi zero.
Stąd $\int_X\mathrm d\eta=0$.

Dla $X=H^k$ identyczny argument działa dla $i<k$.
W ostatniej zmiennej otrzymujemy
\[
 \int_0^\infty \partial_k a_k(u',t)
 \,\mathrm d\lambda_1(t)=-a_k(u',0),
 \qquad u'=(u_1,\ldots,u_{k-1}),
\]
ponieważ $a_k$ znika dla dużych $t$.
Dla $k\ge2$ wynika stąd
\[
 \int_{H^k}\mathrm d\eta
 =-\int_{\R^{k-1}}a_k(u',0)\,\mathrm d\lambda_{k-1}(u').
\]
Na brzegu cofnięcie $\mathrm du_k$ jest zerowe,
więc pozostaje jedynie składnik
$(-1)^{k-1}a_k(u',0)\,
\mathrm du_1\wedge\cdots\wedge\mathrm du_{k-1}$.
Znak orientacji brzegu w tych parametrach
jest równy $(-1)^k$ według
\ref{ma2:prop:boundary-orientation}.
Iloczyn obu znaków wynosi $-1$, co daje
dokładnie tę samą całkę.
Dla $k=1$ obie strony wynoszą $-a_1(0)$,
zgodnie z definicją całki w wymiarze zero.
\end{proof}

\section{Lokalizacja na podrozmaitości}

Do przejścia od modeli lokalnych do ogólnej
podrozmaitości wystarczy skończona lokalizacja
w otoczeniu zwartego nośnika formy.
Nie potrzebujemy tutaj pełnego twierdzenia
o rozkładzie jedności dla dowolnego pokrycia.

\begin{lemat}[Gładka lokalizacja przy zwartym zbiorze]
\label{ma2:lem:stokes-localization}
Niech $M\subseteq\R^n$ będzie gładką
podrozmaitością z brzegiem lub bez brzegu,
a $K\subseteq M$ zbiorem zwartym.
Niech $(V_\alpha)$ będzie względnie otwartym
pokryciem $K$ przez fragmenty parametryzowane.
Istnieją skończenie liczne funkcje gładkie
$\rho_i:M\to[0,1]$ o zwartych nośnikach
zawartych w wybranych zbiorach $V_{\alpha_i}$,
takie że $\sum_i\rho_i=1$ w pewnym względnie
otwartym otoczeniu $K$ w $M$.
\end{lemat}
\begin{proof}
Dla $K=\varnothing$ bierzemy pustą rodzinę.
W przeciwnym razie dla każdego $p\in K$
wybieramy fragment $V_{\alpha(p)}$ zawierający $p$.
Jego dziedzina parametrów jest względnie otwarta
w $\R^k$ lub $H^k$. Wybieramy w niej względnie
otwarte otoczenie punktu odpowiadającego $p$,
o zwartym domknięciu zawartym w tej dziedzinie.
Jego obraz $W_p$ jest względnie otwarty w $M$,
a jego domknięcie w $M$ jest zwarte i zawarte
w $V_{\alpha(p)}$: obraz zwartego domknięcia
w parametrach jest zwarty, a więc domknięty w $M$.

Istnieje $r_p>0$ takie, że
$M\cap B(p,2r_p)\subseteq W_p$, ponieważ
$W_p$ jest względnie otwartym otoczeniem $p$.
Wybieramy gładką funkcję $\psi_p:\R^n\to[0,1]$
równą $1$ na $B(p,r_p)$ i znikającą
poza domkniętą kulą $B(p,2r_p)$.
Dla jawnej konstrukcji niech
\[
 b(t)=
 \begin{cases}\exp(-1/t),&t>0,\\0,&t\le0,\end{cases}
 \qquad
 q(t)=\frac{b(t)}{b(t)+b(1-t)}.
\]
Funkcja $b$ jest gładka: każda pochodna
dla $t>0$ jest iloczynem $\exp(-1/t)$
i wielomianu w $1/t$ (indukcja po rzędzie),
więc dąży do zera przy $t\to0^+$.
Różniczkowanie w punkcie zero daje kolejno
pochodne równe zero, gdyż wykładniczy czynnik
maleje szybciej od dowolnej potęgi $t$.
Mianownik funkcji $q$ jest dodatni,
a $q$ jest równa $0$ dla $t\le0$
i $1$ dla $t\ge1$.
Możemy zatem przyjąć
\[
 \psi_p(x)=q\left(\frac{4r_p^2-\|x-p\|^2}{3r_p^2}\right).
\]
Nośnik ograniczenia $\psi_p$ do $M$
jest zawarty w zwartym domknięciu $W_p$,
a więc w $V_{\alpha(p)}$.

Zwartość $K$ daje skończone podpokrycie
kulami $B(p_i,r_{p_i})$. Niech
$S=\sum_i\psi_{p_i}|_M$.
Na $K$ mamy $S\ge1$.
Definiujemy funkcję gładką na $\R$ przez
\[
 g(t)=
 \begin{cases}q(4t-1)/t,&t>0,\\0,&t\le0.\end{cases}
\]
Jest gładka, gdyż licznik znika dla $t\le1/4$.
Dla $t\ge1/2$ mamy $g(t)=1/t$.
Kładziemy $\rho_i=\psi_{p_i}|_M\,g(S)$.
Nośnik $\rho_i$ zawiera się w nośniku
$\psi_{p_i}|_M$, a wartości należą do $[0,1]$,
bo dla $S>0$ mamy $\psi_{p_i}|_M\le S$.
Na względnie otwartym zbiorze $\{S>1/2\}$,
zawierającym $K$, suma $\rho_i$ jest równa $1$.
\end{proof}

\section{Ogólne twierdzenie Stokesa}

\begin{twierdzenie}[Ogólne twierdzenie Stokesa]
\label{ma2:thm:general-stokes}
Niech $M\subseteq\R^n$ będzie zorientowaną
gładką podrozmaitością wymiaru $k\ge1$
z brzegiem lub bez brzegu, a $\omega$
formą stopnia $k-1$ klasy $C^1$
o zwartym nośniku w $M$.
Brzeg $\partial M$ wyposażamy w orientację
indukowaną przez $M$.
Wtedy
\[
 \boxed{\displaystyle
 \int_M\mathrm d\omega=\int_{\partial M}j^*\omega,}
\]
gdzie $j:\partial M\to M$ jest odwzorowaniem
włączenia.
W szczególności gdy $M$ jest zwarta,
nie trzeba osobno zakładać zwartości nośnika.
\end{twierdzenie}
\begin{proof}
Nośnik $\mathrm d\omega$ zawiera się
w zwartym zbiorze $K=\operatorname{supp}\omega$.
Brzeg jest domknięty w $M$, ponieważ
wnętrze podrozmaitości jest względnie otwarte.
Nośnik $j^*\omega$ zawiera się w $K\cap\partial M$,
który jest zwarty. Obie całki istnieją zatem
na mocy definicji \ref{ma2:def:integral-top-form}
i jej zerowymiarowego odpowiednika.

Pokrywamy $K$ fragmentami parametryzowanymi
$\sigma_i:U_i\to V_i$ takimi, że dziedziny
$U_i$ są spójne i względnie otwarte w modelach
$X_i=\R^k$ lub $H^k$. Możemy to zrobić, ograniczając
mapy do odpowiednio małych kul lub półkul.
Niech $s_i\in\{-1,1\}$ będzie znakiem
orientacji parametryzacji $\sigma_i$.
Z lematu \ref{ma2:lem:stokes-localization}
otrzymujemy skończenie wiele funkcji $\rho_i$,
których suma wynosi $1$ na otoczeniu $K$
i których nośniki są zwarte i zawarte w $V_i$.
Mamy na całym $M$
\[
 \omega=\sum_i\rho_i\omega,\qquad
 \mathrm d\omega=\sum_i\mathrm d(\rho_i\omega),
 \qquad j^*\omega=\sum_i j^*(\rho_i\omega).
\]
Pierwsza równość zachodzi na $K$ dzięki sumie
funkcji, a poza $K$ obie strony są zerowe.
Pozostałe wynikają z liniowości pochodnej
zewnętrznej i cofnięcia.

Ustalmy $i$ i połóżmy $\eta_i=\sigma_i^*(\rho_i\omega)$.
Jest to forma klasy $C^1$ o zwartym nośniku
zawartym w $U_i$: przeciwobraz nośnika
przez homeomorfizm $\sigma_i$ jest zwarty.
Przedłużamy $\eta_i$ przez zero na $X_i$.
Otrzymujemy formę klasy $C^1$, ponieważ
poza jej zwartym nośnikiem jest ona lokalnie
zerowa również przy brzegu $U_i$ w $X_i$.
Pochodna zewnętrzna tego przedłużenia
jest przedłużeniem $\mathrm d\eta_i$ przez zero.

Z przemienności pochodnej zewnętrznej z cofnięciem
(\ref{ma2:thm:exterior-derivative}) i wzoru
na całkowanie w parametrach
(\ref{ma2:thm:integral-top-form-chart}) wynika
\[
 \int_M\mathrm d(\rho_i\omega)
 =s_i\int_{X_i}\mathrm d\eta_i.
\]
Jeżeli $X_i=\R^k$, fragment $V_i$ nie przecina
brzegu $M$, więc wkład $\int_{\partial M}
j^*(\rho_i\omega)$ jest zerowy, podobnie jak
prawa strona z lematu lokalnego.

Jeżeli $X_i=H^k$, ograniczenie $\sigma_i$
do brzegu modelu parametryzuje $V_i\cap\partial M$.
Względem orientacji brzegu modelu ta parametryzacja
ma znak $s_i$, gdyż znak obu orientacji
w standardowych parametrach brzegu wynosi
odpowiednio $(-1)^k$ i $s_i(-1)^k$
(\ref{ma2:prop:boundary-orientation}).
Z własności składania cofnięć
\ref{ma2:prop:form-pullback} i wzoru na całkę
w parametrach dostajemy
\[
 \int_{\partial M}j^*(\rho_i\omega)
 =s_i\int_{\partial X_i}j_i^*\eta_i,
\]
gdzie $j_i:\partial X_i\to X_i$ jest włączeniem.
W wymiarze $k=1$ ta sama równość wynika
bezpośrednio z porównania znaków punktu brzegu.
Lemat \ref{ma2:lem:stokes-local} daje
równość obu wkładów w każdym przypadku.
Sumowanie po skończonej rodzinie $i$ kończy dowód.
\end{proof}

\begin{wniosek}
\label{ma2:cor:stokes-no-boundary}
Jeżeli $\partial M=\varnothing$ i $\omega$
spełnia założenia twierdzenia
\ref{ma2:thm:general-stokes}, to
$\int_M\mathrm d\omega=0$.
\end{wniosek}
\begin{proof}
Całka po pustym brzegu jest zerowa.
\end{proof}

\begin{przyklad}[Twierdzenie Newtona--Leibniza]
\label{ma2:ex:stokes-fundamental}
Niech $M=[a,b]$, gdzie $a<b$, z orientacją
rosnącej współrzędnej, a $f$ będzie funkcją
klasy $C^1$ na $M$. Traktujemy $f$ jako formę
stopnia zero. Mamy $\mathrm df=f'(t)\,\mathrm dt$,
natomiast brzeg składa się z punktu $a$
ze znakiem $-1$ i punktu $b$ ze znakiem $+1$.
Twierdzenie \ref{ma2:thm:general-stokes}
przyjmuje postać
\[
 \int_a^b f'(t)\,\mathrm d\lambda_1(t)=f(b)-f(a).
\]
Nie jest to nowy niezależny dowód twierdzenia
Newtona--Leibniza: użyliśmy go już w dowodzie
lematu \ref{ma2:lem:stokes-local}. Równość pokazuje
natomiast, że ogólna formuła odtwarza podstawową
zasadę rachunku całkowego w wymiarze pierwszym.
\end{przyklad}
